%------------------------------------------------------------------------------%
\section{Equações exponenciais}

As equações em que a incógnita aparece nos expoentes são chamadas de
\emph{Equações exponenciais}\index{Equação!exponencial}.
A seguir serão apresentados alguns exemplos de
equações exponenciais, juntamente com as resoluções.
Para resolvê-las usaremos o fato de que a função exponencial é injetora,
isto é:
\[
  a^{x_1} = a^{x_2}
  \quad\Rightarrow\quad
  x_1 = x_2.
\]

%------------------------------------------------------------------------------%
\begin{example}
  Resolva a equação exponencial \(\left( \frac{1}{3} \right)^x = 81\).

  \solution
  A ideia para resolver uma equação exponencial desse tipo, transformamos
  a equação numa igualdade de potências de mesma base
  \[
    \left( \frac{1}{3}\right)^x = 81
    \quad\Rightarrow\quad
    3^{-x} = 3^4
  \]
  Como a função exponencial é injetora então basta igualar os expoentes,
  o que implica que
  \[
    -x = 4
    \quad\Rightarrow\quad
    x = -4.
  \]
  Logo,
  \[
    S = \{ -4 \}.
  \]
\end{example}
%------------------------------------------------------------------------------%

%------------------------------------------------------------------------------%
\begin{example}
  Resolva a equação exponencial \(5^x-125=0\).

  \solution
  Podemos reescrever a equação como
  \[
    5^x = 5^3,
  \]
  ou seja, transformando em uma igualdade de potências de mesma base.
  Como a função exponencial é injetora então basta igualar os expoentes,
  o que implica que $x=3$.
  Logo,
  \[
    S = \{ 3 \}.
  \]
\end{example}
%------------------------------------------------------------------------------%

%------------------------------------------------------------------------------%
\begin{example}
  Resolva a equação exponencial
  \(
    \left( \frac{1}{2} \right)^x = \sqrt[3]{4}
  \).

  \solution
  Note que
  \begin{align*}
    \left( \frac{1}{2} \right)^x & = \sqrt[3]{4}                      \\[2mm]
    2^{-x}                       & = \left(2^2\right)^{\sfrac{1}{3}}  \\[2mm]
    2^{-x}                       & = 2^{\sfrac{2}{3}}                 \\[2mm]
    -x                           & = \frac{2}{3}                      \\[2mm]
    x                            & =-\frac{2}{3}
  \end{align*}
  Logo,
  \[
    S = \left\{ \frac{-2}{3} \right\}
  \]
\end{example}
%------------------------------------------------------------------------------%

%------------------------------------------------------------------------------%
\begin{example}
  Resolva a equação exponencial \(2^{x-3}+2^{x-1}+2^x=52\).

  \solution
  Nesse caso, não conseguimos transformar a equação em uma igualdade de
  mesma base. Desse modo, vamos reescrever a equação, usando as
  propriedades de potenciação
  \[
    2^x \, 2^{-3} + 2^x \, 2^{-1} + 2^x = 52.
  \]
  Colocando $2^x$ em evidência
  \begin{align*}
    2^x \, 2^{-3} + 2^x \, 2^{-1} + 2^x              & = 52  \\[2mm]
    2^x \left( \frac{1}{8}+  \frac{1}{2} + 1 \right) & = 52  \\[2mm]
    2^x \, \frac{13}{8}                              & = 52  \\[2mm]
    2^x                                              & = 32  \\[2mm]
    2^x                                              & = 2^5 \\[2mm]
    x                                                & = 5
  \end{align*}
  Usamos que 32 pode ser escrito na base 2 como $2^x=2^5$
  Logo, a solução da equação é
  \[
    S = \{ 5 \}.
  \]
\end{example}
%------------------------------------------------------------------------------%

%------------------------------------------------------------------------------%
\begin{example}
  Resolva a equação exponencial \(2^{2x}-9.2^x+8=0\).

  \solution
  Do mesmo modo do item anterior, não conseguimos transformar a equação
  em uma igualdade de mesma base. Vamos reescrever a equação da seguinte
  forma
  \[
    \left(2^x\right)^2 - 9 \cdot 2^x + 8 = 0
  \]
  Agora, vamos fazer a seguinte mudança de variável $y=2^x$. Assim,
  \[
    y^2 - 9y + 8 = 0
  \]
  Que é uma equação do segundo grau, cujas as raízes são $y=1$ e $y=8$.

  Agora, precisamos encontrar os valores de $x$ correspondentes.

  Se $y=1 \Rightarrow 2^x=1 \Rightarrow 2^x=2^0 \Rightarrow x=0$.

  Se $y=8 \Rightarrow 2^x=2^3 \Rightarrow x=3$.

  Logo o conjunto solução da equação é
  \[
     S=\{ 0, 3 \}.
  \]
\end{example}
%------------------------------------------------------------------------------%
